\documentclass[10pt]{article}

% Specify the margins and text size.
\setlength{\textwidth}{6.4in}
\setlength{\textheight}{9.5in}
\setlength{\oddsidemargin}{0pt}
\setlength{\evensidemargin}{0pt}
\setlength{\topmargin}{0pt}
\setlength{\hoffset}{.05in}
\setlength{\voffset}{-.8in}

\setlength{\parskip}{5pt}
\setlength{\parindent}{0pt}

% Load some fonts and symbol packages
\usepackage{latexsym}
\usepackage{pifont}       % contains 'star' symbol for counterinsurgency handbook title
\usepackage{yfonts} 
\usepackage{amsmath}
\usepackage{amsfonts}

\usepackage{graphicx}     % actually, this is loaded by pstricks
\usepackage[T1]{fontenc}
\usepackage{ifthen}
\usepackage{pstricks,pst-grad,pst-text,pst-node,multido,pst-plot,calc,pst-3dplot}
%\usepackage[all]{xy}
%\usepackage{animate}

% The hyperref package inserts links into the pdf.
\definecolor{MyLinkColor}{rgb}{.1,.2,1}
\definecolor{MyCiteColor}{rgb}{.1,1,.2}
\definecolor{MyURLColor}{rgb}{.4,.4,.4}
\usepackage[backref=true,pagebackref=false,hyperindex,colorlinks=true,
  linkcolor=MyLinkColor,urlcolor=MyURLColor]{hyperref}


% The tweaklist package is something I found on the web.  It provides a simple interface
% for making changes to spacing used in the itemize and enumerate environments.  Comment
% this out if you don't care to use tweaklist.
\usepackage{tweaklist}
\renewcommand{\itemhook}{\setlength{\parskip}{2pt}\setlength{\parsep}%
{1pt}\setlength{\topsep}{0pt}\setlength{\itemsep}{0pt}}

\newcommand{\U}{\underline{\hspace{5pt}}}

\usepackage{listings}
\newcounter{EX}\setcounter{EX}{1}
\newcommand{\EXERCISE}{\arabic{EX}.\stepcounter{EX} }

\begin{document}
\pagestyle{empty}
\lstset{language=R, showspaces=false, showstringspaces=false}

\href{http://www.su.edu}{\includegraphics[height=1.75cm]{sulogo.eps}}
\vspace{-1.69cm}

{\small \hfill
\begin{tabular}{cl}
& Math 207\\& Introduction to Statistics\\
\end{tabular}
}
\setlength{\baselineskip}{1.05\baselineskip}
\medskip

\begin{center}
\textbf{\large  Introductory \texttt{R} Exercises}
\end{center}

\newcommand{\SUBX}{\smallskip\hspace{10pt}}
\newcommand{\BSK}{\vspace{.14in}}

\EXERCISE \textbf{Arithmetic Operations}.
Please enter the following commands in \texttt{R} then write down the output
that \texttt{R} produces. Explain to yourself what \texttt{R} thinks it
is doing.
%\vspace{-5pt}

\SUBX a) \texttt{2.0 + 3.1}
\BSK

\SUBX b) \texttt{2.0 * 3.1}
\BSK

\SUBX c) \texttt{3.0 $\widehat{\hspace{3pt}}$ 2.0}
\BSK

\SUBX d) \texttt{9.0  $\widehat{\hspace{3pt}}$  0.5}
\BSK

\SUBX e) \texttt{sqrt(9.0)}
\BSK

\SUBX f) \texttt{exp(1)}
\BSK

\SUBX g) \texttt{sin(pi)}
\BSK

\SUBX h) \texttt{cos(pi/2)}
\BSK

\SUBX i) \texttt{tan(pi/2)}
\BSK

\EXERCISE \textbf{Vector Operations}.
Try the following in \texttt{R} and explain to yourself
what \texttt{R} thinks it is doing.
\vspace{-5pt}

\SUBX a) \texttt{x <- c(1, 0, -2)}
\BSK

\SUBX b) \texttt{x}
\BSK

\SUBX c) \texttt{y <- c(-1, 4.5, 3)}
\BSK

\SUBX d) \texttt{x + y}
\BSK

\SUBX e) \texttt{x + 2}
\BSK

\SUBX f) \texttt{x*y}
\BSK

\SUBX g) \texttt{x  $\widehat{\hspace{3pt}}$  y}
\BSK

\SUBX h) \texttt{exp(x)}
\BSK

\SUBX i) \texttt{1:10}
\BSK

\SUBX j) \texttt{seq(1, 21, by=5)}
\BSK

\vfill
\eject

\EXERCISE \textbf{Statistics}.
Try the following in \texttt{R} and explain to yourself
what \texttt{R} thinks it is doing.

\SUBX a) \texttt{x <- seq(-5, 5)} %\texttt{-5:5!
\BSK

\SUBX b) \texttt{mean(x)}
\BSK

\SUBX c) \texttt{sd(x)}
\BSK

\SUBX d) \texttt{var(x)}
\BSK

\SUBX e) \texttt{min(x)}
\BSK

\SUBX f) \texttt{max(x)}
\BSK

\SUBX g) \texttt{range(x)}
\BSK

\SUBX h) \texttt{length(x)}
\BSK

\SUBX i) \texttt{median(x)}
\BSK

\SUBX j) \texttt{sum(x)}
\BSK

\SUBX k) \texttt{prod(x)}
\BSK

\SUBX $\ell$) Repeat a)-- k) using \texttt{x <- square(seq(0, 2, by=0.2))}.  
Compare your results.
\vspace{.6in}

\EXERCISE \textbf{Subsets}.  Try the following in \texttt{R} and explain to yourself
what \texttt{R} thinks it is doing.

\SUBX a) \texttt{a <- seq(10, 100, by = 10)}
\BSK

\SUBX b) \texttt{a}
\BSK

\SUBX c) \texttt{a[1]}
\BSK

\SUBX d) \texttt{a[5]}
\BSK

\SUBX e) \texttt{a[0]}
\BSK

\SUBX f) \texttt{a[c(2, 4, 6, 8)]}
\BSK

\SUBX g) \texttt{a[-c(2, 4, 6, 8)]}
\BSK

\SUBX h) \texttt{a[2:3] <- c(-333, -555)}
\BSK
\vfill
\eject

\EXERCISE \textbf{Booleans}.  Try the following in \texttt{R} and explain to yourself
what \texttt{R} thinks it is doing.

\SUBX a) \texttt{a <- seq(10, 100, by = 10)}
\BSK

\SUBX b) \texttt{a > 30}
\BSK

\SUBX c) \texttt{a[a > 30]}
\BSK

\SUBX d) \texttt{a == 50}
\BSK

\SUBX e) \texttt{a[a == 50]}
\BSK

\SUBX f) \texttt{a[a != 50]}
\BSK

\SUBX g) \texttt{a <= 30}
\BSK

\SUBX h) \texttt{a[a <= 30]}
\BSK

\SUBX i) \texttt{a[(a >= 50) \& (a <= 90)]}
\BSK

\SUBX j) \texttt{a[(a < 30) | (a > 70)]}
\BSK

\SUBX k) \texttt{a[(a <= 30) \& !(a < 20)]}
\BSK

\EXERCISE \textbf{Missing Values}. Try the following in \texttt{R} and explain to yourself
what \texttt{R} thinks it is doing.

\SUBX a) \texttt{x <- c(NA, 2, 3)}
\BSK

\SUBX b) \texttt{x}
\BSK

\SUBX c) \texttt{x + 8}
\BSK

\SUBX d) \texttt{x < 5}
\BSK

\SUBX e) \texttt{is.na(x)}
\BSK

\SUBX f) \texttt{mean(x)}
\BSK

\SUBX g) \texttt{mean(x, na.rm = TRUE)}
\BSK

\SUBX h) \texttt{x[is.na(x)] = 0}
\BSK

\SUBX i) \texttt{x}
\BSK

\SUBX j) \texttt{mean(x)}
\BSK
\vfill
\eject

\EXERCISE \textbf{Experiments}

\SUBX a) Enter \texttt{sample(10, 4)} several times.  
What is \texttt{R} doing?  
\BSK

\SUBX b) Repeat a) using different values in place of 10 and 4.  What happens
  if you enter \texttt{sample(10, 11)}?
\BSK

\SUBX c) Enter \texttt{sample(c("H", "T"), 5, replace=TRUE)}.
  What is \texttt{R} doing?
\BSK

\SUBX d) Use \texttt{R} to simulate flipping a coin 100 times.  1000 times. 
\bigskip

\EXERCISE \textbf{Graphics}

\SUBX a) Enter the following and sketch the result.

\hspace{20pt} \texttt{x <- seq(-4,4,by=0.1)}

\hspace{20pt} 
\texttt{plot(x, dnorm(x), type="l", cex.lab=1.5, cex.axis=1.5, xlim=c(-4,4), ylim=c(0,0.5))}
\vspace{0.5in}

\SUBX b) Explain what the \texttt{xlim} and \texttt{ylim} parts of the 
  command above mean.
\BSK\BSK

\SUBX c) Experiment to try to figure out what the \texttt{cex.lab}
and \texttt{cex.axis} parts of the command do.
\bigskip

\EXERCISE \textbf{Graphics II}.
Do the following in \texttt{R}.

\begin{verbatim}
install.packages("ggplot2")  # You only need to install this once.
library(ggplot2)             # Do this every time you want to use the package.

?mpg
head(mpg)
summary(mpg)


p <- ggplot(mpg, aes(x = displ, y = hwy))
p <- p + geom_point() # add a plot layer with points
print(p)

p <- ggplot(mpg, aes(x = displ, y = hwy))
p <- p + geom_point(aes(colour = class))
print(p)

p <- ggplot(mpg, aes(x = displ, y = hwy))
p <- p + geom_point(aes(colour = class, size = cyl, shape = drv))
print(p)
\end{verbatim}

\vfill
\eject

\end{document}

